\section{Original arithmetic and what the uniform estimate leaves to prove}
The source is the preceding infinite-array tranche and its pinned finite
SplitZero application \cite{Prior,Boundary}. We retain all original labels,
not only the eigenvalues of their observation. Let $p=4m+1$ be prime. The
canonical row consists of
\[
 m+1\le a\le2m,\quad R=4a-p,\quad
 a=\prod_iq_i^{e_i},\quad -e_i\le\beta_i\le e_i,\quad c\in\{E,M\}.
 \tag{1}
\]
Set
\[
 u=\prod_iq_i^{e_i+\beta_i},\quad g_E=4u+1,\quad g_M=u+a,\quad
 D_\alpha=\frac R{\gcd(R,g_c)},\quad x_\alpha=D_\alpha^{-1}.
 \tag{2}
\]
Every original prime-power defect and channel remains part of $\alpha$.
The exact possible observation values are
\[
 \Sigma=\{0\}\cup\{1/(2k+1):k\ge0\},
 \qquad x_\alpha\in\Sigma\setminus\{0\}.
 \tag{3}
\]
Zero is a possible spectral limit, not an absent labelled state.
The inherited bound treated this set as the larger interval union
$[0,1/3]\cup\{1\}$. The present improvement uses the gaps inside that interval.

For completeness, the positive rational candidate is
\[
 E:\ \left(a,\frac{pa+a^2/u}{R},\frac{pa+p^2u}{R}\right),\qquad
 M:\ \left(a,\frac{p(a+a^2/u)}R,\frac{p(a+u)}R\right).
 \tag{4}
\]
The two nonfirst reduced denominators are both exactly $D_\alpha$.
Indeed $(RY-pa)(RZ-pa)=p^2a^2$, and $4a\equiv p\pmod R$ with
$\gcd(pa,R)=1$ reduces the numerator gcds to those in (2).
Consequently a positive integral witness is exactly a state with $D_\alpha=1$.
These are classical Type I/II coordinates, with the marking retained
\cite[\S2, Propositions 2.2 and 2.6]{ET}.

Writing $d_0=\gcd(a,u)$ gives the complete normalization
\[
 (h,r,s)=(d_0^2/u,u/d_0,a/d_0),\quad a=hrs,\quad u=hr^2,\quad\gcd(r,s)=1.
\]
At a hit the quotient and raw ordered denominators are
\[
 E:\ \kappa=(pr+s)/R,\quad(a,hs\kappa,phr\kappa),\qquad
 M:\ \lambda=(r+s)/R,\quad(a,phs\lambda,phr\lambda).
 \tag{5}
\]
The ordered inverse recovers $u=a^2/(RY-pa)$ for E and
$u=pa^2/(RY-pa)$ for M, then recovers every $\beta_i=v_{q_i}(u)-e_i$.
There is no unmarked interchange of the last two denominators.

\paragraph{Why the row is complete.}
Here is the previously established first-half argument, included to identify
its exact hypothesis. Order a positive integral solution as $x\le y\le z$.
Then $p/4<x\le3p/4<p$. If $x>p/2$, the reciprocal equation implies
$y<p$. Thus $z=pk$ by clearing denominators modulo $p$. Put
$t=(4k-1)/p\in\mathbb Z_{>0}$; it has $t\equiv3\pmod4$,
$\gcd(t,k)=1$, and $txy=k(x+y)$. Write $x=du,y=dv$ with
$\gcd(u,v)=1$, $u\le v$. Then $uv\mid k$, so $k=uvl$ and
$td=l(u+v)$. Coprimality gives $u+v=th,d=lh$.
The inequality $2x>p$ now gives $2lu(u-v)>-1$, forcing $u=v=1$.
But then $t\mid2$, contrary to $t\equiv3\pmod4$.
Equality $x=p/2$ is impossible for odd $p$.
For $a=x$, both other denominators exceed $p/2$, and the positive integer
factor pair $(RY-pa,RZ-pa)$ has product $p^2a^2$. Its two $p$-adic orders are
$(0,2)$, $(1,1)$, or $(2,0)$, yielding exactly (4), with a recorded exterior
reversal in the last case. Thus (1) decides existence without restricting a ray
or a Koide chamber.

Let $W_p$ and $H_p$ be the canonical row's total and successful counts.
The original counting Hilbert space and projector are
\[
 \mathcal H=\ell^2\!\left(\bigsqcup_p\mathscr A_p\right),\quad
 Xe_\alpha=x_\alpha e_\alpha,\quad
 \Pi e_\alpha=\mathbf1_{D_\alpha=1}e_\alpha.
 \tag{6}
\]
The preceding theorem gives $\|X^n-\Pi\|\le3^{-n}$ and compensates row
mass by an explicit prime-dependent power. It does not itself prove $H_p>0$.
The following construction yields one-sided inequalities with faster
convergence and a controlled moment-coefficient norm.

\section{Discrete interpolation filters and their exact constants}
For $j,s\ge0$ define
\[
 P_j(x)=\prod_{\nu=1}^j\frac{(2\nu+1)x-1}{2\nu},\qquad
 F_{j,s}(x)=x^sP_j(x),\qquad \Phi_j=F_{j,1}.
 \tag{7}
\]
These are explicit rational polynomials. In particular $\Phi_j$ is the unique
polynomial of degree at most $j+1$ with values one at $1$ and zero at
$0,1/3,\ldots,1/(2j+1)$. The full evaluation map on these $j+2$ distinct
nodes is an isomorphism; its inverse is
\[
 (v_i)_i\longmapsto\sum_i v_i\prod_{h\ne i}\frac{x-z_h}{z_i-z_h}.
 \tag{8}
\]
This is the classical Lagrange interpolation map \cite[\S3.3(i)]{DLMF}.
Only its stated inverse column is used in (7); no novelty of interpolation
itself is claimed.

\begin{theorem}[Exact one-sided filters]\label{thm:filters}
For $D=2k+1>1$,
\[
 F_{j,s}(1/D)=
 \begin{cases}
 0,&1\le k\le j,\\
 (-1)^j\displaystyle\frac{\binom{k-1}{j}}{(2k+1)^{j+s}},&k>j.
 \end{cases}
 \tag{9}
\]
Also $F_{j,s}(1)=1$, and its coefficient $\ell^1$ norm is exactly $j+1$.
Thus even $j$ gives an upper and odd $j$ a lower bound for the indicator of $D=1$.
On the entire set (3),
\[
 \boxed{\Phi_{2r+1}\le\mathbf1_{\{1\}}\le\Phi_{2r},\qquad
 \|\Phi_j-\mathbf1_{\{1\}}\|_\Sigma
 \le\frac1{2^j j!(2j+3)}.}
 \tag{10}
\]
The even sequence decreases and the odd sequence increases to that indicator.
\end{theorem}
\begin{proof}
Each factor at $D=2k+1$ is $(\nu-k)/(\nu D)$, proving (9) and the sign.
For $k>j$,
$\binom{k-1}{j}/D^j\le1/(2^j j!)$ and $D\ge2j+3$.
This proves (10), including the value zero at the accumulation point.
All nonzero coefficients of $P_j$ alternate in sign, so their absolute sum is
$|P_j(-1)|=\prod_{\nu=1}^j(\nu+1)/\nu=j+1$; multiplication by $x^s$ only
shifts them. For the same-parity comparison, the additional two factors
have product between zero and one at every remaining failed node, and are
one at the successful node. Nodes newly encountered become zero. These
statements prove monotonicity and convergence.
\end{proof}

The sign assertion is only on the actual discrete spectrum. For example
$P_2(1/4)<0$; using it as an upper bound on all of $[0,1/3]$ would be false.
Even return denominators are excluded by the arithmetic proof, not by fiat.

\begin{theorem}[A finite formula for the sharp hybrid error]\label{thm:sharp}
For $j\ge1,s\ge1$, the exact maximum error of $F_{j,s}$ is obtained among
\[
 k\in\mathbb Z,\quad
 \max\{j+1,\lfloor1+j(j+2)/(2s)\rfloor\}
 \le k\le
 \max\{j+1,\lceil j+j(j+2)/(2s)\rceil\}.
 \tag{11}
\]
Evaluate the positive rational expression in (9) at these finitely many
integers and take its maximum. For $j=0,s>0$, the maximum is $3^{-s}$.
For $s=0$ the supremum is $1/(2^j j!)$, approached as $D\to\infty$.
\end{theorem}
\begin{proof}
For real $k>j$, logarithmic differentiation of
$\binom{k-1}{j}/(2k+1)^{j+s}$ gives the sign of
\[
 \sum_{\nu=1}^j\frac{2\nu+1}{k-\nu}-2s.
\]
The sum is strictly decreasing from infinity to zero. Its unique crossing is
between $1+j(j+2)/(2s)$ and $j+j(j+2)/(2s)$, since
$\sum(2\nu+1)=j(j+2)$. Before that crossing the function increases and after it
decreases. The maximizing integer, also respecting $k\ge j+1$, lies in (11).
The boundary cases follow directly from (9).
\end{proof}

In particular the first errors for $\Phi_j$, $j=0,1,2,3$, are
$1/3,2/49,10/2197,4/9261$, attained at $D=3,7,13,21$.
The theorem is a maximization on the full allowed odd spectrum; it does not
assert that each maximizing $D$ occurs in every arithmetic row.

\section{Uniform row positivity, original mass and certified precision}
For any finite original selection $A$, put
\[
 M_k(A)=\sum_{\alpha\in A}D_\alpha^{-k},\qquad
 Q_j(A)=\sum_{\alpha\in A}\Phi_j(D_\alpha^{-1}).
\]
The count $H_A$ retains repeated states with different original labels.
Theorem~\ref{thm:filters} proves
\[
 \boxed{Q_{2r+1}(A)\le H_A\le Q_{2r}(A),\qquad
 |Q_j(A)-H_A|\le\frac{|A|}{2^j j!(2j+3)}.}
 \tag{12}
\]
The same formulas allow positive prescribed masses, but integer recovery below
uses the original unit masses, not a row normalized to total mass one.

The original row bound is
\[
 W_p\le B_p=
 \min\{6m^2,4m(1+\lfloor\log_2(2m)\rfloor)^2\}.
 \tag{13}
\]
Indeed $\tau(a^2)\le2a-1$ gives the first expression. Prime by prime,
$2e+1\le(e+1)(e+2)/2$, so $\tau(a^2)\le d_3(a)$.
Then $\sum_{a\le N}d_3(a)\le N(\sum_{u\le N}1/u)^2$ and dyadic grouping
bounds the harmonic sum by $1+\lfloor\log_2N\rfloor$. This proves (13).

\begin{theorem}[Uniform unscaled lower certificates]\label{thm:rows}
For $n\ge0$, let $j(p,n)$ be the least even nonnegative integer satisfying
\[
 2^j j!(2j+3)\ge8\cdot3^n B_p.
 \tag{14}
\]
Set $\ell_{p,n}=Q_{j+1}(\mathscr A_p)$ and
$u_{p,n}=Q_j(\mathscr A_p)$. Simultaneously for every row,
\[
 \boxed{H_p-\frac1{16\cdot3^n}\le\ell_{p,n}\le H_p
 \le u_{p,n}\le H_p+\frac1{8\cdot3^n}.}
 \tag{15}
\]
In particular $\lceil\ell_{p,n}\rceil=\lfloor u_{p,n}\rfloor=H_p$, and
\[
 \boxed{H_p>0\iff\ell_{p,n}>0.}
 \tag{16}
\]
The largest moment index is $j(p,n)+2$. For fixed $n$, it is
$(1+o(1))\log p/\log\log p$ as $p\to\infty$ through these rows.
\end{theorem}
\begin{proof}
Apply (12)--(14). The ratio of consecutive denominators is
$2(j+1)(2j+5)/(2j+3)>2$, proving the lower error in (15).
The bracket has width at most $3/(16\cdot3^n)<1$ and contains the integer $H_p$.
For $H_p\ge1$ its lower endpoint is positive; for $H_p=0$ it is nonpositive.
All degree choices depend only on the proved mass bound, not on success.
Finally the integral bounds for $\sum_{h=1}^j\log h$ give
$\log(2^j j!(2j+3))=j\log j+O(j)$. Minimality among even integers changes
this by only $O(\log j)$. Since $\log B_p=\log p+O(\log\log p)$,
the asserted degree estimate follows.
\end{proof}

This is a certificate-degree improvement, not a bound on the cost of computing
the original moments. For a mass bound $10^{12}$ the stated schedule needs
moments through degree 14; for $10^{24}$, through degree 22. At the first hard
prime of Section~\ref{sec:hard}, its whole-row mass bound is
$22\,829\,467\,930\,080$, giving maximum moment degree 16. Its whole row was
not enumerated in this contribution.

\paragraph{Absolute input error.}
If every supplied moment through degree $j+2$ has certified absolute error
at most $\delta=(16(j+2)3^n)^{-1}$, polynomial evaluation error is at most
$(j+1)\delta$ for the upper filter and $(j+2)\delta$ for the lower.
These follow from the exact coefficient norm in Theorem~\ref{thm:filters}.
Enlarge both endpoints in the required directions. The resulting interval
still contains $H_p$ and has width less than $7/(16\cdot3^n)<1$.
Original masses are unchanged; if moments are large, obtaining this absolute
precision still requires the corresponding relative precision.

\paragraph{Infinite operators.}
On the original countable space (6), the pointwise inequalities give
$\Phi_{2r+1}(X)\preceq\Pi\preceq\Phi_{2r}(X)$ with the same operator-norm errors.
Let $N e_p=H_pe_p$ on the maximal domain
$\{v:\sum H_p^2|v_p|^2<\infty\}$. The row operators with diagonals
$\ell_{p,n},u_{p,n}$ have that same domain; their differences from $N$ are
bounded with the constants in (15). This follows directly by coordinate norms,
so no unproved trace of an infinite identity is used. For real $\sigma>1$,
the signed relative row traces converge absolutely, bounded by the corresponding
constant in (15) times $\sum_p p^{-\sigma}$.

For the whole programme, (16) gives
\[
 \forall p\ H_p>0
 \iff\inf_p\ell_{p,0}>0
 \iff\inf_p\ell_{p,0}\ge15/16.
 \tag{17}
\]
This is a precise forcing target, not a proof of the last bound. An empty row
would have every odd-filter sum nonpositive. The present argument does not
rule out such a row merely by assembling all rows.

\section{Source-resolved lower bounds and exact transport}
For a genuine partition $A=\bigsqcup_\lambda A_\lambda$, retain its original
unit masses. Let $L_d(\mu_\lambda)$ be the previous signed variational bound
\[
 L_d(\mu_\lambda)=\max_{\deg P\le d,\ P(1)=1}
 \frac32\int(x-1/3)P(x)^2d\mu_\lambda.
 \tag{18}
\]
Its continuous-support sign is valid on $[0,1/3]\cup\{1\}$; the discrete
polynomials (7) use the stronger exact spectrum and belong to a different
family. We do not claim one family dominates the other at every degree.

\begin{proposition}[A lower bound retaining the partition]
The integer
\[
 C_d(\{A_\lambda\})=
 \sum_\lambda\max\{0,\lceil L_d(\mu_\lambda)\rceil\}
 \tag{19}
\]
is a lower bound for $H_A$. Refining the original partition cannot decrease it.
In particular a positive score on one original channel certifies an ES witness
without requiring a positive sum after merging the channels.
\end{proposition}
\begin{proof}
Each $L_d(\mu_\lambda)\le H_{A_\lambda}$ and the latter is a nonnegative
integer. For a refinement of one part, every common polynomial has score the
sum of its refined scores, so $L_{\rm coarse}\le\sum L_{\rm fine}$.
Use $\lceil\sum x_i\rceil\le\sum\lceil x_i\rceil$ and discard negative
summands to obtain (19) and its monotonicity.
\end{proof}

At degree one, the two-channel coefficient space has block form
$\operatorname{diag}(J_E,J_M)$. Restriction through
$\Delta P=(P,P)$ has form $J_E+J_M$. This is restriction to a particular
subspace, not an equivalence of the forms. The inverse of $\Delta$ on its
image reads either identical entry; the difference component must be retained
for the full direct sum.

\paragraph{Split arithmetic.}
For a budget split $e=f+b$, the addition map has fibre
\[
 \max(-f_i,\beta_i-b_i)\le\xi_i\le\min(f_i,\beta_i+b_i),\quad\eta=\beta-\xi,
\]
with cardinality $n(\beta)$ equal to the product of those interval lengths.
Assign each occurrence weight $1/n(\beta)$. For every new filter and every
moment, the occurrence sum equals the original sum. This is proved by grouping
the actual finite fibres, not by assuming an unweighted convolution.
In the corresponding completed spaces, duplication $U$ and fibre averaging
$E=U^*$ satisfy $EU=1$. Since $\Phi_j(0)=0$,
\[
 \boxed{\Phi_j(UXU^*)=U\Phi_j(X)U^*.}\tag{20}
\]
Its inverse on $\operatorname{im}U$ uses $U^*$, and the complementary kernel
is retained. The raw repeated diagonal operator still counts split
multiplicities; (20) is its compressed transport, not a licence to take the raw
trace. For $P_j$ without the extra $x$, the missing term would be
$(-1)^j(2^j j!)^{-1}(I-UU^*)$.

Every polynomial observation retains its source state $\alpha$ and $D_\alpha$.
Forgetting them can merge zeroed defects $3,5,\ldots,2j+1$. A scalar moment
vector does not reconstruct which divisor supplied a hit. Positive lower
scores imply that some original summand is positive; enumerating its finite
original fibre then returns a $D=1$ label and (5). The prime-power socle and
augmentation maps commute with these blockwise scalar polynomials, without
turning supported zeros into external absence.

\section{An actual hard-prime obstruction and two positive repairs}\label{sec:hard}
Take
\[
 \boxed{\begin{gathered}p_*=19\,748\,674\,681\equiv1\pmod{840},\qquad R=135,\\
 a_*=4\,937\,168\,704=2^6\cdot13^4\cdot37\cdot73.\end{gathered}}
 \tag{21}
\]
This is a prime, certified by the full order criterion with base $11$ and
\[
 p_*-1=2^3\cdot3\cdot5\cdot7\cdot23510327.
\]
All factors are proved prime by complete trial division. The residues of
$11^{(p_*-1)/q}$, in this factor order, are
\[
 19748674680,\ 10039250736,\ 3428156451,\ 13907516733,\ 12674160518.
\]
Each displayed residue minus one is coprime to $p_*$, and
$11^{p_*-1}=1\pmod{p_*}$.
For any prime divisor $l$ of $p_*$, these identities force the order modulo
$l$ to be divisible by $p_*-1$, hence $l\ge p_*$. This proves primality.
The program separately tests every possible odd divisor through $140529$.

The following table exhausts the $1053$ divisors and both original channels:
\[
\begin{array}{c|rrrrrrrr}
D&1&3&5&9&15&27&45&135\\\hline
E&1&39&53&80&69&144&244&423\\
M&0&40&54&80&68&144&244&423\\
E+M&1&79&107&160&137&288&488&846
\end{array}\tag{22}
\]
It can be independently reproduced in the finite group algebra by multiplying
\[
 (\sum_{i=0}^{12}[2]^i)(\sum_{j=0}^8[13]^j)
 (\sum_{k=0}^2[37]^k)(\sum_{l=0}^2[73]^l)
 \quad\text{in }\mathbb Z[(\mathbb Z/135\mathbb Z)^\times]
 \tag{23}
\]
and applying the two gcd gates to each residue. The independent direct
exponent enumeration is also supplied, with all 2106 labels. No group-generated
support is substituted for these bounded exponents.

For $S_i=\sum_\alpha\gcd(R,g_\alpha)^i$, the first values are
\[
(S_0,S_1,S_2,S_3,S_4)
=(2106,13962,315738,12455226,722183418).
\]
Let $A_i=3S_{i+1}-RS_i$. Then
\[
 (A_0,A_1,A_2)=(-242424,-937656,-5258952),\qquad
 A_0A_2-A_1^2=395697405312>0.
 \tag{24}
\]
The degree-one signed localizing matrix is congruent, by a positive diagonal
coordinate change, to $\begin{psmallmatrix}A_0&A_1\\A_1&A_2\end{psmallmatrix}$.
It is negative definite. Thus the shell has a hit but \emph{every} aggregate
degree-one normalized polynomial fails to give a positive score. Its exact
maximum is $-50887012/144801015$.
This is an obstruction to extending the earlier finite observation, not a
counterexample to ES or a correction of the proved sufficient criterion.
No leastness is claimed.

\paragraph{First repair: retain the channel.}
On E alone, the fixed polynomial $P(x)=(36x-1)/35$ has score
\[
 \boxed{\frac32\sum_{\alpha\in E}(x_\alpha-1/3)P(x_\alpha)^2
 =\frac{64376}{196875}>0.}\tag{25}
\]
The exact E maximum is $1692013/5172810>0$; the M maximum is negative.
Thus (19) forces at least one original E witness.

\paragraph{Second repair: exploit the discrete spectrum.}
The fixed quartic filter
\[
 \Phi_3(x)=\frac{105x^4-71x^3+15x^2-x}{48}
\]
gives, on the original \emph{combined} state set,
\[
 \boxed{Q_3=\frac{1297691}{2460375}>0,
 \qquad Q_4=\frac{115374406}{110716875}.}\tag{26}
\]
Their integer bracket is exactly $H=1$. In particular, the arithmetic
inequality
\[
 \boxed{105S_4-71RS_3+15R^2S_2-R^3S_1>0}\tag{27}
\]
is a valid sufficient condition for a hit in \emph{every} original shell.
It is satisfied here although (24) prevents first-degree certification.
This uses one higher moment for positivity, with no singular-Gram inverse.

The unique original hit in this shell is
\[
 (c,u,h,r,s,\kappa)=
 (E,141879296,26,2336,81289,341725215823),
\]
\[
 (x,y,z)=
 (4937168704,722241027794932022,409884191461517484986392768).
 \tag{28}
\]
The centered vector is $(5,-3,-1,1)$. Its slope $2336/81289<1/8$ also belongs
to the prescribed lower positive chamber. That chamber assertion is only for
this marking, not every later member of the progression below.

\paragraph{Channel separation alone is not asserted universal.}
There is a second fully certified hard prime
\[
 p_{**}=2671180768668258904496300170662649\equiv289\pmod{840},
\]
with $a=(p_{**}+135)/4=2^5\cdot431^3\cdot457^3\cdot7229^4$.
Its 9702 canonical candidates contain three E hits, but
\[
 L_1(E+M)=-\frac{7138170316}{1971871105},\qquad
 L_1(E)=-\frac{146243172}{493377905},\qquad L_1(M)<0.
\]
The fixed discrete quartic nevertheless gives
\[
 Q_3=\frac{54448414}{66430125}>0.
\]
The exact full histogram, three witnesses and recursive Lucas certificate
are in the JSON. The eight nontrivial certificate nodes reduce to trial-prime
leaves below 10000; no probable-prime assumption is used. This additional case
excludes a premature claim that channel separation always suffices.

\section{A constructive infinite arithmetic specialization}
The first recovered witness extends pointwise, without requiring its parameter
to be prime. For every integer $j\ge0$, put
\[
 t=73+945j,\quad C=67632448,\quad n=4Ct-135
 =19748674681+255650653440j,
\]
\[
 h_0=26,\quad r_0=32t,\quad s_0=81289,\quad a=Ct,
\]
\[
 \kappa_0=341725215823+8847406287328j+57265746370560j^2.
\]
Polynomial expansion proves $135\kappa_0=nr_0+s_0$ and
$a=h_0r_0s_0$. Therefore
\[
 \boxed{\frac4n=\frac1a+\frac1{h_0s_0\kappa_0}
                  +\frac1{nh_0r_0\kappa_0}.}\tag{29}
\]
All denominators are positive integers for every $j\ge0$.
If $d=\gcd(r_0,s_0)$, then $\gcd(d,135)=1$ and $d\mid\kappa_0$.
The canonical return is
\[
 (h,r,s,\kappa)=(h_0d^2,r_0/d,s_0/d,\kappa_0/d),\quad
 u=h_0r_0^2,
\]
which leaves the ordered denominators unchanged. The actual inverse parameter
is $j=(n-19748674681)/255650653440$ on this image.
Every $n$ in (29) is $1\pmod{840}$, and the initial term is coprime to the
step. Dirichlet's theorem therefore gives infinitely many prime members
\cite[Theorem 18.1]{Dirichlet}. Formula (29) already proves ES for every integer
member. This is a specialization of the inherited selector, not a claim that
arithmetic-progression constructions are new or that this progression exhausts
the hard locus.

\section{Verification, source scope and next pointwise obligation}
The standalone deterministic standard-library verifier proves the two primes by exact
certificates, enumerates both specified shells in full, compares direct and
residue-algebra censuses, verifies the moment inequalities and their positive
witness inverses, and tests the entire finite node window used for sharp errors.
It also checks interpolation inverse columns, interval precision, completed
split trace identities and a full small-prime row comparison. Its finite output
is not a proof of the infinite estimates; those were proved above.

The original infinite-array and SplitZero source measures, kernels and labels
are retained. No theta-source norm, Frobenius action, Leech geometry or cubic
surface existence theorem is imported to assert positivity of these arithmetic
sums. The generic interpolation and moment principles have their primary
citations. The application and examples were calculated here; a scoped search
did not establish historical priority.

The uniform estimates now give explicit original-moment separating inequalities,
with factorial rather than merely geometric error and controlled coefficient
size. Two actual hard-prime shells show why a fixed first-degree shortcut cannot
be inferred from the earlier census, and how stronger signed estimates recover
positive witnesses. What is not proved is that every canonical row defined here,
for primes $p\equiv1\pmod4$, must satisfy one
of these positive lower inequalities. The remaining universal assertion is the
arithmetic lower bound in (17), not an unproved limit interchange.
